\section{Agent Security Sandbox}
\label{sec:framework}

We introduce the \textbf{Agent Security Sandbox} (\asb{}), a modular framework for evaluating LLM agent defenses against indirect prompt injection.
\asb{} consists of four components: a threat model, a benchmark suite, a defense-strategy library, and an automated evaluation pipeline.

\subsection{Threat Model}
\label{sec:threat_model}

We consider an LLM agent $\mathcal{A}$ that receives a user goal $g$, invokes tools $T = \{t_1, \ldots, t_k\}$, and retrieves untrusted external content $c$ (e.g., email bodies, search results, file contents) during execution.

\textbf{Attacker.}
The adversary controls $c$ and embeds an injection $\iota$ designed to cause $\mathcal{A}$ to execute a forbidden action $a^* \notin \text{Auth}(g)$, where $\text{Auth}(g)$ is the set of actions authorized by the user's goal.
The adversary has \emph{black-box} access: no knowledge of system prompts, model weights, or conversation history.
In the adaptive setting (\S\ref{sec:adaptive}), the adversary additionally observes defense rejection signals and iteratively refines $\iota$.

\textbf{Defender.}
The defender may modify the prompt (prompt-level), intercept tool calls (tool-gating), or filter outputs (content-level), but cannot modify the base model.
The defender knows the agent's tool set $T$ and the user's goal $g$, but not the specific injection $\iota$.

\textbf{Success criteria.}
An attack succeeds if $\mathcal{A}$ executes any action matching a forbidden action pattern $a^*$, regardless of whether other legitimate actions are also performed.
A defense succeeds on benign cases if $\mathcal{A}$ executes all expected tool calls specified by the ground-truth case.

\subsection{Benchmark Suite}
\label{sec:benchmark}

\paragraph{Design Principles.}
The benchmark is designed around four principles:
(1)~\emph{Diversity}: attacks span multiple types, techniques, locations, and tools;
(2)~\emph{Realism}: attack payloads use techniques observed in real-world injection attempts;
(3)~\emph{Balance}: benign cases cover diverse task categories to accurately measure false positives;
(4)~\emph{Reproducibility}: all cases are deterministic with rule-based evaluation.

\paragraph{Attack Cases (352).}
\asb{} includes 352 attack cases spanning 6 types: \emph{data exfiltration} (188, 53.4\%), \emph{goal hijacking} (49, 13.9\%), \emph{privilege escalation} (41, 11.6\%), \emph{multistep} (34, 9.7\%), \emph{denial of service} (20, 5.7\%), and \emph{social engineering} (20, 5.7\%).
Cases employ 54 injection techniques (direct override, fake system messages, encoding-based evasion, RAG poisoning, etc.) placed across 8 locations (email bodies 43.8\%, search results 16.8\%, file contents 16.8\%, and others).

\paragraph{Benign Cases (213).}
Benign cases cover 7 categories including advanced multi-step tasks (60), multilingual (44), untrusted-but-non-malicious content (40), basic single-tool (20), multi-tool workflows (20), edge cases (15), and generated diverse tasks (14).
Each case specifies a user goal $g$, untrusted content $c$, expected tool calls, forbidden actions, and metadata (attack type, technique, location, difficulty).

\begin{table}[t]
\centering
\small
\caption{ASB benchmark statistics. The benchmark comprises 565 cases spanning 6 attack types, 11 tools, and 3 difficulty levels.}
\label{tab:benchmark_stats}
\begin{tabular}{lrrr}
\toprule
\textbf{Attack Type} & \textbf{Easy} & \textbf{Medium} & \textbf{Hard} \\
\midrule
Data Exfiltration & 58 & 36 & 94 \\
Goal Hijacking & 13 & 19 & 17 \\
Privilege Escalation & 5 & 9 & 27 \\
Multi-step & 0 & 3 & 31 \\
Denial of Service & 5 & 7 & 8 \\
Social Engineering & 5 & 11 & 4 \\
\midrule
\textbf{Total Attack} & 86 & 85 & 181 \\
\midrule
Benign & 40 & 125 & 48 \\
\midrule
\textbf{Grand Total} & 126 & 210 & 229 \\
\bottomrule
\end{tabular}
\end{table}


\subsection{Defense Taxonomy}
\label{sec:taxonomy}

Table~\ref{tab:taxonomy} presents the 11 defense strategies in \asb{}, organized by mechanism type.
Strategies D0--D7 are implementations or adaptations of published approaches; D8 and D9 extend ideas from the security literature; D10 (\civ{}) is our novel contribution (details in Appendix~\ref{sec:civ}).

\begin{table*}[t]
\centering
\small
\caption{Defense taxonomy in \asb{}. Mechanism types: \textbf{P}rompt-level, \textbf{T}ool-gating, \textbf{C}ontent-level, \textbf{M}ulti-signal. D10 (\civ{}) is our novel contribution.}
\label{tab:taxonomy}
\begin{tabular}{llllccc}
\toprule
ID & Name & Type & Key Mechanism & Modifies & Gates & Filters \\
   &      &      &               & Prompt   & Tools & Output \\
\midrule
D0  & Baseline           & ---  & No defense                          & --- & --- & --- \\
D1  & Spotlighting        & P    & Delimiter-based source marking      & \checkmark & & \\
D5  & Sandwich            & P    & Goal-reminder wrapping              & \checkmark & & \\
D7  & Input Classifier    & P    & Injection pattern removal           & \checkmark & & \\
\midrule
D2  & Policy Gate         & T    & Risk-level + whitelist enforcement  & & \checkmark & \\
D3  & Task Alignment      & T    & LLM goal--action consistency check  & & \checkmark & \\
D4  & Re-execution        & T    & Clean re-run comparison             & & \checkmark & \\
D8  & Semantic Firewall   & T    & Embedding-based drift detection     & & \checkmark & \\
D9  & Dual-LLM            & T    & Two-model screening                 & & \checkmark & \\
\midrule
D6  & Output Filter       & C    & Regex-based data leak detection     & & & \checkmark \\
\midrule
D10 & \textbf{CIV (ours)} & M    & Provenance + embed.\ compat.\ + plan dev. & & \checkmark & \\
\bottomrule
\end{tabular}
\end{table*}


\paragraph{Prompt-Level Defenses.}
\defense{D1} (Spotlighting) wraps untrusted content in delimiters with explicit warnings \citep{hines2024defending}.
\defense{D5} (Sandwich) re-anchors the model to the original goal after presenting untrusted content \citep{yi2023benchmarking,hines2024defending}.
\defense{D7} (Input Classifier) pre-processes inputs to detect and remove injection patterns.

\paragraph{Tool-Gating Defenses.}
\defense{D2} (Policy Gate) enforces tool-call permissions based on risk levels and whitelists.
\defense{D3} (Task Alignment) uses LLM-based reasoning to verify goal--action consistency.
\defense{D4} (Re-execution) re-runs the agent on sanitized input and compares outputs \citep{xiang2024melon}.
\defense{D8} (Semantic Firewall) uses embedding similarity to detect goal--action semantic drift.
\defense{D9} (Dual-LLM) employs a separate screening model to verify proposed actions \citep{chen2024struq}.

\paragraph{Content-Level Defense.}
\defense{D6} (Output Filter) applies post-hoc regex-based detection of sensitive data leakage.

\paragraph{Multi-Signal Defense.}
\defense{D10} (\civ{}) fuses three verification signals---entity provenance tracking, embedding-based semantic compatibility, and plan deviation detection---to determine whether each tool call is consistent with the user's original goal.
It applies contextual integrity theory \citep{nissenbaum2004privacy} to detect information-flow violations in agent behavior.
Details are in Appendix~\ref{sec:civ}.

\subsection{Evaluation Pipeline}
\label{sec:evaluation}

For each (model, defense, run) triple, \asb{} executes the agent on all benchmark cases with the defense active, then applies a rule-based judge to determine verdicts.
We report three metrics:
\begin{itemize}
    \item \textbf{\asr{}} (Attack Success Rate): fraction of attacks that bypassed the defense and executed the forbidden action. Lower is better.
    \item \textbf{\bsr{}} (Benign Success Rate): fraction of legitimate tasks completed successfully. Higher is better.
    \item \textbf{\fpr{}} (False Positive Rate): fraction of benign tasks incorrectly blocked. $\text{FPR} = 1 - \text{BSR}$.
\end{itemize}
All metrics are reported with 95\% Wilson score confidence intervals \citep{wilson1927probable}.
Pairwise defense comparisons use McNemar's test \citep{mcnemar1947note} with Bonferroni correction.

\asb{} is designed for extensibility: new defenses implement a two-method interface, new cases follow a JSONL schema, and new models are supported through an OpenAI-compatible API.
All experiments are parameterized through YAML configuration files.
